%% Les deux chaînes fonctionnelles d'un système et l'ACTION (Tle, p. 244-246).
%% Chaîne d'information (bleu) au-dessus, chaîne d'énergie (orange) en dessous ;
%% l'ordre issu du traitement arrive sur DISTRIBUER (et non sur ALIMENTER).
\documentclass[tikz,border=4pt]{standalone}
\usepackage{liaisons}
\begin{document}
\begin{tikzpicture}[x=1cm,y=1cm,
  bloc/.style={line width=0.9pt, fill=white, rounded corners=2pt, minimum width=2cm,
               minimum height=0.7cm, inner sep=1pt, font=\scriptsize},
  bi/.style={bloc, draw=solideB},
  be/.style={bloc, draw=solideD},
  fi/.style={-{Stealth[length=5pt,width=4pt]}, line width=1pt, draw=solideB},
  fe/.style={-{Stealth[length=5pt,width=4pt]}, line width=1pt, draw=solideD},
  bandeau/.style={rounded corners=6pt, dash pattern=on 3pt off 2pt, line width=0.8pt},
  et/.style={font=\scriptsize, inner sep=2pt, align=center,
             execute at begin node={\hyphenpenalty=10000\exhyphenpenalty=10000}}]

%% ================== bandeau haut : chaîne d'information ===============
\draw[bandeau, draw=solideB!55] (0.0,3.55) rectangle (6.7,4.95);
\node[font=\scriptsize\bfseries, text=solideB, anchor=south west, inner sep=1pt]
     at (0.0,4.98) {Chaîne d'information};
\node[bi] (acq) at (1.1,4.25) {ACQUÉRIR};
\node[bi] (tra) at (3.35,4.25) {TRAITER};
\node[bi] (com) at (5.6,4.25) {COMMUNIQUER};
\draw[fi] (acq) -- (tra);
\draw[fi] (tra) -- (com);

%% ================== bandeau bas : chaîne d'énergie ====================
\draw[bandeau, draw=solideD!55] (0.0,0.95) rectangle (11.2,2.35);
\node[font=\scriptsize\bfseries, text=solideD, anchor=north west, inner sep=1pt]
     at (0.0,0.92) {Chaîne d'énergie};
\node[be] (ali) at (1.1,1.65) {ALIMENTER};
\node[be] (dis) at (3.35,1.65) {DISTRIBUER};
\node[be] (con) at (5.6,1.65) {CONVERTIR};
\node[be] (tr2) at (7.85,1.65) {TRANSMETTRE};
\node[be] (agi) at (10.1,1.65) {AGIR};
\foreach \a/\b in {ali/dis,dis/con,con/tr2,tr2/agi} { \draw[fe] (\a) -- (\b); }

%% ================== ordres : l'information commande DISTRIBUER ========
\draw[-{Stealth[length=7pt,width=7pt]}, line width=2.4pt, solideD] (3.35,3.5) -- (3.35,2.05);
\node[et, text=solideD, anchor=west] at (3.47,2.8) {\textbf{Ordres}};

%% ================== entrées et sorties de la chaîne d'information =====
\draw[fi] (-1.0,4.25) -- (acq.west);
\node[et, anchor=south east, align=right, text width=1.7cm] at (0.05,4.33) {Consignes, commandes};
\draw[fi] (com.east) -- (8.0,4.25);
\node[et, anchor=south] at (7.3,4.32) {Messages};

%% ================== retour de l'ACTION vers « Informations » ==========
\draw[fi, rounded corners=4pt] (12.3,3.05) -- (12.95,3.05) -- (12.95,5.55) -- (1.1,5.55) -- (1.1,4.98);
\node[et, anchor=south west, inner sep=1pt] at (1.15,5.6) {Informations};

%% ================== zone ACTION =======================================
\draw[draw=solideA, line width=1.1pt, fill=solideA!8] (11.4,0.75) rectangle (12.3,3.4);
\node[font=\small\bfseries, text=solideA, rotate=90] at (11.85,2.07) {ACTION};
\draw[fe] (agi.east) -- (11.4,1.65);
%% grosses flèches creuses : matière d'œuvre entrante puis sortante
\draw[draw=solideA, line width=0.9pt, fill=solideA!15]
  (11.63,4.55) -- (12.07,4.55) -- (12.07,3.85) -- (12.25,3.85) -- (11.85,3.42) -- (11.45,3.85)
  -- (11.63,3.85) -- cycle;
\node[et, text=solideA, anchor=south, text width=2cm] at (11.85,4.62) {Matière d'œuvre entrante};
\draw[draw=solideA, line width=0.9pt, fill=solideA!15]
  (11.63,0.72) -- (12.07,0.72) -- (12.07,0.1) -- (12.25,0.1) -- (11.85,-0.33) -- (11.45,0.1)
  -- (11.63,0.1) -- cycle;
\node[et, text=solideA, anchor=north, text width=2cm] at (11.85,-0.4) {Matière d'œuvre sortante};

%% ================== pertes ============================================
\draw[fe] (10.1,1.28) -- (10.75,0.62);
\node[et, text=solideD, anchor=north, inner sep=1pt] at (10.72,0.56) {Pertes};
\end{tikzpicture}
\end{document}
